\chapter*{\centering PERNYATAAN PENGGUNAAN \textit{AI-GENERATED CONTENT}}
\addcontentsline{toc}{chapter}{PERNYATAAN PENGGUNAAN \textit{AI-GENERATED CONTENT}}
\begin{spacing}{1.15}

\noindent Kami yang bertanda tangan di bawah ini, pengusul proposal Tugas Akhir yang berjudul:
\begin{center}
  \textbf{[Judul Usulan Penelitian]}
\end{center}
dengan ini menyatakan bahwa:
\begin{enumerate}
  \item Kami bertanggung jawab penuh atas seluruh isi, keakuratan, dan orisinalitas tulisan dalam proposal ini.
  \item Pemanfaatan teknologi \textit{Artificial Intelligence} (AI) sebagai alat bantu penulisan tidak menggantikan proses berpikir kritis, analisis, interpretasi, dan penyusunan argumen ilmiah yang kami lakukan.
  \item Penggunaan AI dalam penyusunan naskah ini kami ungkapkan secara jujur dan bertanggung jawab sebagaimana tercantum pada tabel berikut.
  \item Naskah ini telah diperiksa menggunakan perangkat deteksi AI (\textit{AI detector}) premium yang diinformasikan oleh Koordinator Tugas Akhir dengan hasil indikasi penggunaan AI sebesar \textbf{[xx]\%}, yaitu tidak melebihi batas 20\%. Bukti hasil pemeriksaan dilampirkan pada Lampiran~1.
\end{enumerate}

% Isi tabel berikut dengan seluruh penggunaan AI. Jika AI tidak digunakan sama
% sekali, ganti tabel dengan kalimat: "Kami tidak menggunakan AI dalam
% penyusunan naskah ini."
\begin{center}
  \begin{tabularx}{\textwidth}{|>{\centering\arraybackslash}p{0.06\textwidth}|*{3}{>{\RaggedRight\arraybackslash}X|}}
  \hline
  \textbf{No.} & \textbf{Alat AI (Versi)} & \textbf{Bagian Naskah} & \textbf{Bentuk Penggunaan} \\ \hline
  1 & [Nama alat, versi] & [Bab/subbab] & [Misalnya: parafrase, pemeriksaan tata bahasa, penerjemahan] \\ \hline
  2 & [Nama alat, versi] & [Bab/subbab] & [Bentuk penggunaan] \\ \hline
  \end{tabularx}
\end{center}

\noindent Apabila di kemudian hari ditemukan bukti bahwa pernyataan ini tidak benar, kami bersedia menerima sanksi sesuai dengan peraturan yang berlaku di Institut Teknologi Del.

\vspace{0.3cm}
\begin{flushright}
  Sitoluama, [tanggal bulan tahun]
\end{flushright}

\noindent
\begin{tabularx}{\textwidth}{>{\centering\arraybackslash}X>{\centering\arraybackslash}X>{\centering\arraybackslash}X}
  Pengusul 1, & Pengusul 2, & Pengusul 3, \\[2cm]
  \underline{Nama Mahasiswa 1} & \underline{Nama Mahasiswa 2} & \underline{Nama Mahasiswa 3} \\
  NIM 1 & NIM 2 & NIM 3 \\
\end{tabularx}
\end{spacing}
